\documentclass[11pt]{article}

\usepackage[margin=1in]{geometry}
\usepackage{amsmath, amssymb, amsthm}
\usepackage{mathtools}
\usepackage{enumitem}
\usepackage{booktabs}
\usepackage{tikz-cd}
\usepackage[colorlinks=true, linkcolor=blue, citecolor=blue]{hyperref}

% ---------- Theorem environments ----------
\theoremstyle{plain}
\newtheorem{theorem}{Theorem}[section]
\newtheorem{proposition}[theorem]{Proposition}
\newtheorem{lemma}[theorem]{Lemma}
\newtheorem{corollary}[theorem]{Corollary}

\theoremstyle{definition}
\newtheorem{definition}[theorem]{Definition}
\newtheorem{example}[theorem]{Example}
\newtheorem{exercise}{Exercise}

\theoremstyle{remark}
\newtheorem*{remark}{Remark}

% ---------- Macros ----------
\newcommand{\Z}{\mathbb{Z}}
\newcommand{\Q}{\mathbb{Q}}
\newcommand{\R}{\mathbb{R}}
\newcommand{\C}{\mathbb{C}}
\newcommand{\gen}[1]{\langle #1 \rangle}
\DeclareMathOperator{\im}{im}
\DeclareMathOperator{\GL}{GL}
\DeclareMathOperator{\SL}{SL}
\DeclareMathOperator{\sgn}{sgn}
\DeclareMathOperator{\ord}{ord}

\title{Introduction to Group Theory\\[0.5ex]
  \large Lecture Notes for Math 3110: Abstract Algebra I}
\author{Department of Mathematics}
\date{Fall 2026}

\begin{document}

\maketitle

\begin{abstract}
These notes cover the first three weeks of the course: the definition of a group, subgroups, cyclic groups, cosets and Lagrange's theorem, normal subgroups and quotient groups, and homomorphisms up to the First Isomorphism Theorem. Proofs are given in full for the main results. Exercises at the end are of the kind you will see on the problem sets. For a fuller treatment, see \cite{DF} or \cite{Artin}.
\end{abstract}

\tableofcontents

\section{Groups}

\begin{definition}\label{def:group}
A \emph{group} is a set $G$ together with a binary operation $G \times G \to G$, written $(a, b) \mapsto ab$, such that
\begin{enumerate}[label=\textup{(G\arabic*)}]
  \item \label{ax:assoc} the operation is associative: $(ab)c = a(bc)$ for all $a, b, c \in G$;
  \item \label{ax:id} there is an element $e \in G$, the \emph{identity}, with $ea = ae = a$ for all $a \in G$;
  \item \label{ax:inv} every $a \in G$ has an \emph{inverse} $a^{-1} \in G$ with $aa^{-1} = a^{-1}a = e$.
\end{enumerate}
The group is \emph{abelian} if in addition $ab = ba$ for all $a, b$. The \emph{order} of $G$, written $|G|$, is the number of its elements.
\end{definition}

For abelian groups we often write the operation as $+$, the identity as $0$ and the inverse of $a$ as $-a$.

\begin{example}\label{ex:groups}
\leavevmode
\begin{enumerate}[label=(\alph*)]
  \item $(\Z, +)$, $(\Q, +)$, $(\R, +)$ and $(\C, +)$ are abelian groups. The integers under multiplication are \emph{not} a group: $2$ has no multiplicative inverse in $\Z$.
  \item $\Q^\times = \Q \setminus \{0\}$ is an abelian group under multiplication, and so are $\R^\times$ and $\C^\times$.
  \item For $n \ge 1$, the set $\Z/n\Z = \{0, 1, \dots, n-1\}$ with addition modulo $n$ is an abelian group of order $n$.
  \item The \emph{general linear group} $\GL_n(\R)$ of invertible $n \times n$ real matrices is a group under matrix multiplication. It is non-abelian for $n \ge 2$.
  \item The \emph{symmetric group} $S_n$ of all bijections $\{1, \dots, n\} \to \{1, \dots, n\}$, under composition, has order $n!$ and is non-abelian for $n \ge 3$.
  \item The \emph{dihedral group} $D_n$ of symmetries of a regular $n$-gon has order $2n$. It is generated by a rotation $r$ through $2\pi/n$ and a reflection $s$, subject to
  \[
    r^n = s^2 = e, \qquad srs = r^{-1}.
  \]
\end{enumerate}
\end{example}

\begin{proposition}\label{prop:basic}
Let $G$ be a group and $a, b, c \in G$.
\begin{enumerate}[label=\textup{(\roman*)}]
  \item The identity of $G$ is unique.
  \item Each element has a unique inverse.
  \item $(ab)^{-1} = b^{-1}a^{-1}$ and $(a^{-1})^{-1} = a$.
  \item \textup{(Cancellation)} If $ab = ac$, then $b = c$; if $ba = ca$, then $b = c$.
\end{enumerate}
\end{proposition}

\begin{proof}
(i) If $e$ and $e'$ are both identities, then $e = ee' = e'$, using first that $e'$ is an identity and then that $e$ is.

(ii) If $b$ and $c$ are both inverses of $a$, then $b = be = b(ac) = (ba)c = ec = c$.

(iii) Using associativity, $(ab)(b^{-1}a^{-1}) = a(bb^{-1})a^{-1} = aa^{-1} = e$, and similarly $(b^{-1}a^{-1})(ab) = e$. By (ii), $b^{-1}a^{-1}$ is \emph{the} inverse of $ab$. The second statement holds because $a$ satisfies the defining property of an inverse of $a^{-1}$.

(iv) Multiply $ab = ac$ on the left by $a^{-1}$ and use associativity. The right-hand version is symmetric.
\end{proof}

A finite group can be described completely by its \emph{Cayley table}, whose entry in row $a$ and column $b$ is the product $ab$. Table~\ref{tab:klein} shows the Klein four-group $V_4 = \{e, a, b, c\}$, in which every non-identity element is its own inverse and the product of any two distinct non-identity elements is the third.

\begin{table}[ht]
  \centering
  \begin{tabular}{c|cccc}
    $\cdot$ & $e$ & $a$ & $b$ & $c$ \\ \hline
    $e$ & $e$ & $a$ & $b$ & $c$ \\
    $a$ & $a$ & $e$ & $c$ & $b$ \\
    $b$ & $b$ & $c$ & $e$ & $a$ \\
    $c$ & $c$ & $b$ & $a$ & $e$
  \end{tabular}
  \caption{Cayley table of the Klein four-group $V_4$.}
  \label{tab:klein}
\end{table}

By cancellation (Proposition~\ref{prop:basic}(iv)), every row and every column of a Cayley table contains each element exactly once.

\section{Subgroups}

\begin{definition}
A subset $H \subseteq G$ is a \emph{subgroup} of $G$, written $H \le G$, if $H$ is a group under the operation of $G$ restricted to $H$.
\end{definition}

Checking all three axioms every time is unnecessary. Associativity is inherited from $G$, and the remaining conditions can be packed into a single test.

\begin{proposition}[Subgroup test]\label{prop:subgroup-test}
A nonempty subset $H$ of a group $G$ is a subgroup if and only if $ab^{-1} \in H$ for all $a, b \in H$.
\end{proposition}

\begin{proof}
If $H \le G$, then $b^{-1} \in H$ and $H$ is closed under the operation, so $ab^{-1} \in H$.

Conversely, suppose the condition holds. Since $H \neq \varnothing$, pick $a \in H$. Then $e = aa^{-1} \in H$. For any $b \in H$, $b^{-1} = eb^{-1} \in H$. Finally, for $a, b \in H$ we now know $b^{-1} \in H$, so $ab = a(b^{-1})^{-1} \in H$. Thus $H$ contains the identity, is closed under inverses and products, and inherits associativity from $G$.
\end{proof}

\begin{example}\label{ex:subgroups}
\leavevmode
\begin{enumerate}[label=(\alph*)]
  \item For each $n \ge 0$, $n\Z = \{nk : k \in \Z\}$ is a subgroup of $\Z$, since $na - nb = n(a - b)$.
  \item $\SL_n(\R) = \{A \in \GL_n(\R) : \det A = 1\}$ is a subgroup of $\GL_n(\R)$, because $\det(AB^{-1}) = \det A / \det B = 1$.
  \item The \emph{center} $Z(G) = \{z \in G : zg = gz \text{ for all } g \in G\}$ is a subgroup of $G$.
\end{enumerate}
\end{example}

\begin{definition}
For $g \in G$, the \emph{cyclic subgroup generated by $g$} is $\gen{g} = \{g^k : k \in \Z\}$. The \emph{order} of $g$, written $\ord(g)$, is the smallest positive integer $n$ with $g^n = e$, or $\infty$ if there is none.
\end{definition}

\begin{proposition}\label{prop:order}
Let $g \in G$ have finite order $n$. Then
\begin{enumerate}[label=\textup{(\roman*)}]
  \item $g^m = e$ if and only if $n \mid m$;
  \item $\gen{g} = \{e, g, g^2, \dots, g^{n-1}\}$, and these $n$ elements are distinct, so $|\gen{g}| = \ord(g)$.
\end{enumerate}
\end{proposition}

\begin{proof}
(i) If $n \mid m$, say $m = qn$, then $g^m = (g^n)^q = e$. Conversely, suppose $g^m = e$ and divide: $m = qn + r$ with $0 \le r < n$. Then
\[
  g^r = g^{m - qn} = g^m (g^n)^{-q} = e.
\]
Since $n$ is the smallest positive exponent giving $e$, we must have $r = 0$.

(ii) Every $g^k$ equals $g^r$ where $r$ is the remainder of $k$ modulo $n$, by the same computation. If $g^i = g^j$ with $0 \le i < j < n$, then $g^{j-i} = e$ with $0 < j - i < n$, which contradicts the minimality of $n$.
\end{proof}

\section{Cyclic groups}

\begin{definition}
A group $G$ is \emph{cyclic} if $G = \gen{g}$ for some $g \in G$, called a \emph{generator}.
\end{definition}

The groups $\Z$ and $\Z/n\Z$ are cyclic, generated by $1$. Every cyclic group is abelian, since $g^i g^j = g^{i+j} = g^j g^i$.

\begin{theorem}\label{thm:cyclic-subgroups}
Every subgroup of a cyclic group is cyclic.
\end{theorem}

\begin{proof}
Let $G = \gen{g}$ and $H \le G$. If $H = \{e\}$, then $H = \gen{e}$. Otherwise $H$ contains some $g^k$ with $k \neq 0$, and then also $g^{-k}$, so $H$ contains a positive power of $g$. Let $m$ be the smallest positive integer with $g^m \in H$. We claim $H = \gen{g^m}$.

Clearly $\gen{g^m} \subseteq H$. Conversely, let $g^k \in H$ and write $k = qm + r$ with $0 \le r < m$. Then
\[
  g^r = g^k (g^m)^{-q} \in H,
\]
because $H$ is a subgroup. By the minimality of $m$, $r = 0$, so $g^k = (g^m)^q \in \gen{g^m}$.
\end{proof}

\begin{corollary}\label{cor:subgroups-of-Z}
The subgroups of $\Z$ are exactly the sets $n\Z$ for $n \ge 0$.
\end{corollary}

\begin{proposition}\label{prop:Zn-generators}
In $\Z/n\Z$, the element $k$ has order $n / \gcd(n, k)$. In particular, $k$ generates $\Z/n\Z$ if and only if $\gcd(n, k) = 1$, so $\Z/n\Z$ has exactly $\varphi(n)$ generators, where $\varphi$ is Euler's totient function.
\end{proposition}

The proof is Exercise~\ref{ex:Zn-order}.

\section{Cosets and Lagrange's theorem}

\begin{definition}
Let $H \le G$ and $g \in G$. The \emph{left coset} of $H$ containing $g$ is $gH = \{gh : h \in H\}$. Right cosets $Hg$ are defined similarly. The number of distinct left cosets is the \emph{index} $[G : H]$.
\end{definition}

\begin{lemma}\label{lem:cosets}
Let $H \le G$ and $a, b, g \in G$.
\begin{enumerate}[label=\textup{(\roman*)}]
  \item $aH = bH$ if and only if $a^{-1}b \in H$.
  \item Two left cosets of $H$ are either equal or disjoint.
  \item The map $H \to gH$, $h \mapsto gh$, is a bijection. In particular $|gH| = |H|$.
\end{enumerate}
\end{lemma}

\begin{proof}
(i) If $aH = bH$, then $b = be \in bH = aH$, so $b = ah$ for some $h \in H$ and $a^{-1}b = h \in H$. Conversely, if $a^{-1}b = h \in H$, then $b = ah$ and $bH = a(hH) = aH$, since $hH = H$ for $h \in H$.

(ii) Suppose $c \in aH \cap bH$, say $c = ah_1 = bh_2$. Then $a^{-1}b = h_1h_2^{-1} \in H$, and $aH = bH$ by (i).

(iii) The map $x \mapsto g^{-1}x$ from $gH$ to $H$ is an inverse.
\end{proof}

\begin{theorem}[Lagrange]\label{thm:lagrange}
Let $G$ be a finite group and $H \le G$. Then $|H|$ divides $|G|$, and
\[
  [G : H] = \frac{|G|}{|H|}.
\]
\end{theorem}

\begin{proof}
Every $g \in G$ lies in the coset $gH$, so the left cosets cover $G$. By Lemma~\ref{lem:cosets}(ii) distinct cosets are disjoint, so they partition $G$. By Lemma~\ref{lem:cosets}(iii) each coset has exactly $|H|$ elements. If there are $k = [G : H]$ cosets, counting the elements of $G$ gives $|G| = k\,|H|$.
\end{proof}

Lagrange's theorem has several immediate consequences.

\begin{corollary}\label{cor:lagrange}
Let $G$ be a finite group and $g \in G$.
\begin{enumerate}[label=\textup{(\alph*)}]
  \item $\ord(g)$ divides $|G|$.
  \item $g^{|G|} = e$.
  \item If $|G| = p$ is prime, then $G$ is cyclic and every non-identity element generates $G$.
\end{enumerate}
\end{corollary}

\begin{proof}
(a) Apply Theorem~\ref{thm:lagrange} to $H = \gen{g}$ and use Proposition~\ref{prop:order}(ii). (b) follows from (a) and Proposition~\ref{prop:order}(i). For (c), let $g \neq e$. Then $\ord(g)$ divides $p$ and is greater than $1$, so it equals $p$ and $\gen{g} = G$.
\end{proof}

\begin{corollary}[Fermat's little theorem]\label{cor:fermat}
If $p$ is prime and $p \nmid a$, then $a^{p-1} \equiv 1 \pmod{p}$.
\end{corollary}

\begin{proof}
The nonzero residues modulo $p$ form a group $(\Z/p\Z)^\times$ of order $p - 1$ under multiplication. Apply Corollary~\ref{cor:lagrange}(b) to the class of $a$.
\end{proof}

\begin{remark}
The converse of Lagrange's theorem is false. The alternating group $A_4$ has order $12$ but no subgroup of order $6$. For finite \emph{cyclic} groups, though, there is exactly one subgroup of each order dividing $|G|$.
\end{remark}

\section{Normal subgroups and quotient groups}

We would like to turn the set of cosets $G/H$ into a group by setting $(aH)(bH) = abH$. This only works for a special class of subgroups.

\begin{definition}
A subgroup $N \le G$ is \emph{normal}, written $N \trianglelefteq G$, if $gNg^{-1} = N$ for every $g \in G$. Equivalently, $gN = Ng$ for every $g \in G$.
\end{definition}

\begin{example}\label{ex:normal}
\leavevmode
\begin{enumerate}[label=(\alph*)]
  \item Every subgroup of an abelian group is normal.
  \item $\SL_n(\R) \trianglelefteq \GL_n(\R)$, because $\det(gAg^{-1}) = \det A$.
  \item The subgroup $H = \gen{(1\,2)} = \{e, (1\,2)\}$ of $S_3$ is not normal: conjugating by $(1\,3)$ gives
  \[
    (1\,3)(1\,2)(1\,3)^{-1} = (2\,3) \notin H.
  \]
\end{enumerate}
\end{example}

\begin{proposition}\label{prop:index-two}
Every subgroup of index $2$ is normal.
\end{proposition}

\begin{proof}
If $[G : N] = 2$, the two left cosets are $N$ and $G \setminus N$, and so are the two right cosets. If $g \in N$, then $gN = N = Ng$. If $g \notin N$, then $gN = G \setminus N = Ng$.
\end{proof}

In particular the alternating group $A_n$, which has index $2$ in $S_n$, is normal.

\begin{theorem}\label{thm:quotient}
Let $N \trianglelefteq G$. Then the set of cosets $G/N = \{gN : g \in G\}$ is a group under
\[
  (aN)(bN) = abN,
\]
with identity $N$ and inverses $(aN)^{-1} = a^{-1}N$. If $G$ is finite, $|G/N| = |G|/|N|$.
\end{theorem}

\begin{proof}
The only issue is whether the product is well defined, that is, independent of the chosen representatives. Suppose $aN = a'N$ and $bN = b'N$. Then $a' = an_1$ and $b' = bn_2$ for some $n_1, n_2 \in N$, and
\[
  a'b' = a n_1 b n_2 = ab\,(b^{-1} n_1 b)\, n_2.
\]
Normality gives $b^{-1}n_1b \in N$, so $a'b' \in abN$ and $a'b'N = abN$. Associativity, the identity and inverses are inherited from $G$. The statement about orders is Theorem~\ref{thm:lagrange}.
\end{proof}

\begin{example}
The group $\Z/n\Z$ is literally the quotient of $\Z$ by its normal subgroup $n\Z$: the coset $k + n\Z$ is the residue class of $k$ modulo $n$.
\end{example}

\section{Homomorphisms}

\begin{definition}
A map $\varphi\colon G \to H$ between groups is a \emph{homomorphism} if $\varphi(ab) = \varphi(a)\varphi(b)$ for all $a, b \in G$. Its \emph{kernel} and \emph{image} are
\[
  \ker\varphi = \{g \in G : \varphi(g) = e_H\}, \qquad \im\varphi = \{\varphi(g) : g \in G\}.
\]
A bijective homomorphism is an \emph{isomorphism}. If one exists, we write $G \cong H$.
\end{definition}

\begin{example}\label{ex:homs}
\leavevmode
\begin{enumerate}[label=(\alph*)]
  \item $\det\colon \GL_n(\R) \to \R^\times$ is a surjective homomorphism with kernel $\SL_n(\R)$.
  \item The sign map $\sgn\colon S_n \to \{\pm 1\}$ is a homomorphism with kernel $A_n$.
  \item $\exp\colon (\R, +) \to (\R_{>0}, \cdot)$ is an isomorphism, with inverse $\log$.
  \item Reduction modulo $n$, $\Z \to \Z/n\Z$, is a surjective homomorphism with kernel $n\Z$.
\end{enumerate}
\end{example}

\begin{proposition}\label{prop:hom-basic}
Let $\varphi\colon G \to H$ be a homomorphism. Then
\begin{enumerate}[label=\textup{(\roman*)}]
  \item $\varphi(e_G) = e_H$ and $\varphi(g^{-1}) = \varphi(g)^{-1}$;
  \item $\ker\varphi \trianglelefteq G$ and $\im\varphi \le H$;
  \item $\varphi$ is injective if and only if $\ker\varphi = \{e_G\}$.
\end{enumerate}
\end{proposition}

\begin{proof}
(i) $\varphi(e) = \varphi(ee) = \varphi(e)\varphi(e)$, and cancelling gives $\varphi(e) = e$. Then $\varphi(g)\varphi(g^{-1}) = \varphi(e) = e$.

(ii) If $k \in \ker\varphi$ and $g \in G$, then $\varphi(gkg^{-1}) = \varphi(g)\,e\,\varphi(g)^{-1} = e$, so $g(\ker\varphi)g^{-1} \subseteq \ker\varphi$; applying this to $g^{-1}$ gives the reverse inclusion. That $\ker\varphi$ and $\im\varphi$ are subgroups follows from (i) and Proposition~\ref{prop:subgroup-test}.

(iii) $\varphi(a) = \varphi(b)$ if and only if $\varphi(a^{-1}b) = e$, that is, $a^{-1}b \in \ker\varphi$. So $\varphi$ is injective exactly when $a^{-1}b \in \ker\varphi$ forces $a = b$, which happens exactly when the kernel is trivial.
\end{proof}

The main structural result of this part of the course says that every homomorphic image of $G$ is a quotient of $G$.

\begin{theorem}[First Isomorphism Theorem]\label{thm:first-iso}
Let $\varphi\colon G \to H$ be a homomorphism with kernel $K$. Then
\[
  \bar\varphi\colon G/K \to \im\varphi, \qquad gK \mapsto \varphi(g),
\]
is a well-defined isomorphism. Equivalently, $\varphi = \bar\varphi \circ \pi$, where $\pi\colon G \to G/K$ is the quotient map, and the following diagram commutes:
\[
\begin{tikzcd}[column sep=large, row sep=large]
  G \arrow[r, "\varphi"] \arrow[d, "\pi"'] & \im\varphi \\
  G/K \arrow[ru, "\bar\varphi"', "\cong"] &
\end{tikzcd}
\]
\end{theorem}

\begin{proof}
For $g, g' \in G$,
\[
  gK = g'K \iff g^{-1}g' \in K \iff \varphi(g)^{-1}\varphi(g') = e \iff \varphi(g) = \varphi(g'),
\]
using Lemma~\ref{lem:cosets}(i) and Proposition~\ref{prop:hom-basic}(i). Read from left to right, this says $\bar\varphi$ is well defined. Read from right to left, it says $\bar\varphi$ is injective. It is surjective onto $\im\varphi$ by construction, and it is a homomorphism because
\[
  \bar\varphi\bigl((gK)(g'K)\bigr) = \bar\varphi(gg'K) = \varphi(gg') = \varphi(g)\varphi(g') = \bar\varphi(gK)\,\bar\varphi(g'K). \qedhere
\]
\end{proof}

\begin{corollary}\label{cor:iso-examples}
\leavevmode
\begin{enumerate}[label=(\alph*)]
  \item $\GL_n(\R)/\SL_n(\R) \cong \R^\times$.
  \item $S_n/A_n \cong \Z/2\Z$ for $n \ge 2$.
  \item $\Z/n\Z$ is, up to isomorphism, the only cyclic group of order $n$, and $\Z$ is the only infinite cyclic group.
\end{enumerate}
\end{corollary}

\begin{proof}
Apply Theorem~\ref{thm:first-iso} to the homomorphisms of Example~\ref{ex:homs}(a) and (b). For (c), if $G = \gen{g}$, the map $\Z \to G$, $k \mapsto g^k$, is a surjective homomorphism whose kernel is $n\Z$ when $\ord(g) = n$ (Proposition~\ref{prop:order}(i)) and $\{0\}$ when $g$ has infinite order.
\end{proof}

\section{Exercises}

\begin{exercise}
Show that a group in which $g^2 = e$ for every $g$ is abelian. \emph{Hint:} expand $(ab)^2 = e$.
\end{exercise}

\begin{exercise}\label{ex:Zn-order}
Prove Proposition~\ref{prop:Zn-generators}. \emph{Hint:} $mk \equiv 0 \pmod n$ if and only if $\frac{n}{\gcd(n,k)}$ divides $m$.
\end{exercise}

\begin{exercise}
Write out the Cayley table of $\Z/4\Z$. Show that $\Z/4\Z \not\cong V_4$ by comparing the orders of their elements.
\end{exercise}

\begin{exercise}
List all subgroups of $S_3$ and determine which are normal. Verify Lagrange's theorem in each case.
\end{exercise}

\begin{exercise}
Let $H, K \le G$. Prove that $H \cap K \le G$. Give an example in which $H \cup K$ is not a subgroup.
\end{exercise}

\begin{exercise}
Use Corollary~\ref{cor:lagrange} to prove Euler's theorem: if $\gcd(a, n) = 1$, then $a^{\varphi(n)} \equiv 1 \pmod n$.
\end{exercise}

\begin{exercise}
Let $\varphi\colon \Z/12\Z \to \Z/12\Z$ be given by $\varphi(x) = 3x$. Find $\ker\varphi$ and $\im\varphi$, and identify the isomorphism given by Theorem~\ref{thm:first-iso}.
\end{exercise}

\begin{thebibliography}{9}

\bibitem{Artin}
M.~Artin, \emph{Algebra}, 2nd ed., Pearson, Boston, 2011.

\bibitem{DF}
D.~S. Dummit and R.~M. Foote, \emph{Abstract Algebra}, 3rd ed., John Wiley \& Sons, Hoboken, NJ, 2004.

\bibitem{Gallian}
J.~A. Gallian, \emph{Contemporary Abstract Algebra}, 9th ed., Cengage Learning, Boston, 2017.

\end{thebibliography}

\end{document}
